\documentclass{article}
\usepackage[utf8]{inputenc}
\usepackage[T1]{fontenc}
\usepackage[margin=2cm]{geometry}
\usepackage{hyperref}
\sloppy
\begin{document}
\section*{Historical source}
Encyclopaedia Britannica scan, source row 1678. Parse 3.5 extracted the text. Super 3.5 VL corrected it from an enlarged crop of each block. It can still contain recognition errors.
Page headers, footers, and page numbers are removed.

\subsection*{Untitled continuation (E0146)}
\textit{Continues an entry from an earlier page.}

Bilfinger.

Billinger. Edward, had already had occasion to appreciate his talents, and put them to use. At the time when he carried on the war in Servia, he maintained a regular correspondence with Billinger who had long been known as an able Engineer, and had, indeed, made some improvements of value in the received system of Fortification. After his return to Tubingen, he had frequent conversations with the Professor, on different subjects of administration and appointed him, in 1735, Privy-Counsellor. This nomination was not a simple honorary title. Billinger saw himself raised at once to a power almost unlimited. He resisted some time a promotion, which he did not think himself qualified to sustain. In accepting office, his first care was to acquire all the knowledge necessary to the discharge of its duties. He employed almost two years in assiduous labour to instruct himself thoroughly in the statistics of the country, in considering its political situation, its constitution, its interests, and became, at the end of all this study, one of the most enlightened ministers that his country had yet produced. Billinger was placed in a situation too elevated not to excite jealousy and hatred. He felt it, and wished to quit the ministry; but the court refused to receive his resignation, soon after the tender of which, the Duke died. Billinger experienced from his successor all the consideration and all the friendship which he had experienced in the beginning of his career. Received into a confidence without bounds, he had the power to realize, without obstacle, those plans of administration which the most enlightened patriotism had inspired him. Wurtemberg still feels the happy influence of his ministry. Commerce, public instruction, agriculture, were protected and ameliorated by his cares. The culture of the vine, of so much importance in this country, was one of the principal objects of his attention. We ought not to forget that he was the original author of that strict union which has long united Wurtemberg and Prussia, and of the importance to which the hereditary Prince of Wurtemberg was raised at the court of Berlin. In 1737, the Duke nominated him President of his consistory, and Secretary of the grand Order of the Chace. He was also Curator of the University of Tubingen, and member of the Royal Academy of

Berlin. All his time was consecrated to some serious occupation, with the exception of one hour in the evening, which he employed in making and receiving visits. His greatest enjoyment was in cultivating his garden. A warm and strong friend, he gave many proofs of gratitude to those protectors, who had generously assisted him in his studies. He has been reproached with being irascible; but in spite of some slight blemishes, the memory of Bilfinger will be always dear to his countrymen, and honoured by all Germans. Wurtemberg reckons him among the greatest men which she has produced, and proposes him as a model to her statesmen and her men of letters. He was never married, and left no issue. He died at Stuttgart the 18th of February 1750. His works, besides various Papers published in the Memoirs of the St Petersburg and Paris Academies of Science, are :

1. Disputatio de Harmoniá praestabilitá, Tubingen, 1721, in 4to.

2. De Harmonia Animi et Corporis Humani maxime praestabilitá Commentatio Hypothetica. Francfort on the Main, 1723, in 8vo. This work was inserted in the Expurgatory Index at Rome in 1784.

3. De Origine et Permissione Mali, precipue Moralis, Commentatio Philosophica, ibid. 1724, in 8vo.

4. Specimen Doctrinæ Veterum Sinarum Moralis et Politicæ, Francf. 1724, in 4to.

5. Dissertatio Historico-catoptrica de Speculo Archi-
medis, Tubingen, 1725, in 4to.

6. Dilucidationes Philosophicae de Deo, Animâ Humana, Mundo, et Generalibus Rerum Affectionibus, ibid. 1725, in 4to.

7. Bilfingeri et Holmanni Epistola de Harmonia Prestabilitá, 1728, in 4to.

8. Disputatio de Naturâ et Legibus Studii in Theologia Thetici, ibid. 1731, in 4to.

9. Disputatio de Cultu Dei Rationali, ibid. 1731.
10. Notae Breves in Ben. Spinosam Methodum Explicandi Scripturas, Tub. 1732, in 4to.

plicandi Scripturas, Tub. 1732, in 4to.
11. De Mysteriis Christianae Fidei Generatim Spec-
tatis Sermo, Recitatus 1732, Tubingen, 1732, in 4to.
12. Elementa Physices, Leipsig. 1749, in 8vo.

12. Elementa Physices, Leipzig, 1742, in 8vo.

12. Elementa Physices, Leipzig, 1742, in 8vo.
13. La Citadelle Coupée, Leipzig, 1756, in 4to.

See Biographie Universelle, Tom. IV. (z.)


\section*{Source record}
\begin{description}
\item[Scan ID] \texttt{britannica:309f765b1ecc83e0a7f7762e44b4a927a7555443:train:1678}
\item[Image file] \texttt{page-1678-75b494c1b4fb.jpg}
\item[Image SHA-256] \texttt{75b494c1b4fb241bd9dcaba505dfb1732b3c54ff9f843f0d4a7c1c31110a6ec8}
\item[Dataset] NationalLibraryOfScotland/encyclopaedia\_britannica\_illustrated at revision \texttt{309f765b1ecc83e0a7f7762e44b4a927a7555443} (CC0-1.0)
\item[Extraction] model \texttt{nvidia/nemotron-3.5-parse}, raw reply \texttt{extract/raw/page-1678-75b494c1b4fb.json}
\item[Raw reply SHA-256] \texttt{a2443ca44edf0f6b8555c6e9929b715f1eea6ca81932e77270791874d7ff39cd}
\item[Correction] model \texttt{nvidia/Nemotron-3.5-Super-VL}, block crops, raw replies \texttt{correct/raw/crop\_nothink/page-1678-75b494c1b4fb-b*.json}; super\_vl: 16
\item[Running header] B I L B I L
\end{description}
\end{document}
